% Problem record op_0ff53a57552ddd7f. This TeX file is derived from
% database/problems_json/op_0ff53a57552ddd7f.json by scripts/sync-tex.mjs; edit the JSON record
% and rerun the script rather than editing this file.
\section{Unrestricted quantum time--space tradeoff for collision finding}
\label{sec:op_0ff53a57552ddd7f}

\paragraph{Problem.}
Does every unrestricted quantum collision-finding algorithm satisfy $T^2S=\Omega(N\log N)$?

Let $f:[N]\to[N]$ be uniformly random and available through a coherent value oracle. An algorithm must output distinct $x,x'$ with $f(x)=f(x')$ with probability at least $2/3$. Let $T$ be its number of oracle queries and $S$ its maximum retained workspace in qubits, including quantum memory and classical data retained for later use.

Is the tradeoff

\begin{equation}
T^2S=\Omega(N\log N)
\label{eq:0ff5-1}
\end{equation}

valid without any symmetry restriction? In Eq.~\eqref{eq:0ff5-1}, the external oracle storage is excluded from $S$, but query registers and accumulated tables are included.

\subsection*{Status}

Unsolved

\subsection*{Source}

The question is explicitly posed or retained as open in the cited primary literature \sourcecite{ref:0ff5-magniez26}{Magniez26}.  The statement is rewritten here to make its hypotheses and success criterion self-contained.

\subsection*{Progress}

\begin{itemize}
  \item BHT upper bound. Storing a table of $r$ input–output pairs and applying quantum search gives

\begin{equation}
T=O\!\left(r+\sqrt{N/r}\right),\qquad S=O(r\log N).
\label{eq:0ff5-3}
\end{equation}

The displayed definitions, constraints, and target bounds are recorded in Eqs.~\eqref{eq:0ff5-3}.

  \item For $1\le r\le N^{1/3}$, this matches the conjectured tradeoff. The original BHT work introduced the search-based collision method and its time–space interpolation. \sourcecite{ref:0ff5-brassard98}{Brassard98}, \sourcecite{ref:0ff5-magniez26}{Magniez26}

  \item 2023. Hamoudi and Magniez established a lower bound $T^3S=\Omega(K^3N)$ for finding $K$ collision pairs. This does not give the stronger single-collision tradeoff stated here. \sourcecite{ref:0ff5-hamoudi23}{Hamoudi23}

  \item September 9, 2026. Magniez and Zur prove the desired bounds for label-symmetric algorithms, whose strategy respects permutations of the function's output labels. Section 1.4 explicitly leaves removal of this assumption open. Standard symmetrization can require $\Theta(N\log N)$ additional storage and therefore does not preserve the space parameter. \sourcecite{ref:0ff5-magniez26}{Magniez26}

  \item Retained as open in the unrestricted model. The newest paper resolves an important restricted case, not the full statement.
\end{itemize}

\subsection*{References}

\begin{enumerate}[label={},leftmargin=4.8em,labelsep=0.5em]
  \footnotesize
  \item[\textup{[Brassard98]}]\label{ref:0ff5-brassard98}
  Gilles Brassard, Peter Høyer, and Alain Tapp. \emph{Quantum Algorithm for the Collision Problem}. LATIN 1998; \href{https://arxiv.org/abs/quant-ph/9705002}{arXiv:quant-ph/9705002} (1997).

  \item[\textup{[Magniez26]}]\label{ref:0ff5-magniez26}
  Frédéric Magniez and Sebastian Zur. \emph{Tight Time-Space Lower Bounds for Collision Finding and Element Distinctness under Label Symmetry}. \href{https://arxiv.org/html/2609.10808v1}{arXiv:2609.10808}, September 9, 2026, preprint. See Section 1.4 and Theorem 4.18.

  \item[\textup{[Hamoudi23]}]\label{ref:0ff5-hamoudi23}
  Yassine Hamoudi and Frédéric Magniez. \emph{Quantum Time-Space Tradeoff for Finding Multiple Collision Pairs}. ACM Transactions on Computation Theory 15(1–2) (2023); \href{https://arxiv.org/abs/2002.08944v5}{arXiv:2002.08944v5}.
\end{enumerate}

\subsection*{Comment}

This problem is an explicit open extension of a September 9, 2026 preprint.

This asks whether a quantum algorithm can obtain the full collision speedup while retaining asymptotically less usable information than the known table-based methods. The target concerns memory as well as queries, rather than a new query lower bound in the unlimited-space model.

The symmetry issue is substantive: symmetry of a problem does not guarantee a cost-free symmetric implementation of every algorithm. A useful lower-bound invariant would have to control retained information without assuming that the algorithm ignores the numerical identities of output labels. A counterexample would instead need to exploit those identities in a way that improves the query–memory product, not merely the internal implementation constant.

\subsection*{Field}

Quantum algorithm

\subsection*{Topic}

Computational complexity and computability

\subsection*{ID}

\texttt{op\_0ff53a57552ddd7f}
